% Application Architecture Diagram (Standalone)
% Shows the relationship between Application Controller, Application Master,
% Callback Module (.app file), and the Supervision Tree.
% Compile with: xelatex application_architecture.tex
\documentclass[border=10pt]{standalone}

% ============================================================
% Packages
% ============================================================
\usepackage{fontspec}
\usepackage{xeCJK}
\setCJKmainfont[BoldFont=STHeiti, ItalicFont=STKaiti]{STSong}
\setCJKsansfont{STHeiti}
\setCJKmonofont{STFangsong}

\usepackage{xcolor}
\usepackage{tikz}
\usepackage{pifont}  % for \ding symbols (circled numbers)

\usetikzlibrary{arrows.meta, positioning, shapes.geometric, fit, calc, backgrounds, decorations.pathreplacing}

% ============================================================
% Color Definitions
% ============================================================
\definecolor{erlred}{HTML}{A4070A}
\definecolor{erlblue}{HTML}{2B5EA7}
\definecolor{erlgreen}{HTML}{2E7D32}
\definecolor{erlyellow}{HTML}{F9A825}
\definecolor{erlpurple}{HTML}{6A1B9A}
\definecolor{erlmodule}{HTML}{D84315}  % dedicated color for callback module
\definecolor{erlinit}{HTML}{00897B}    % dedicated vivid teal color
\definecolor{erlpid}{HTML}{E65100}     % dedicated color for Pid return
\definecolor{erlapp}{HTML}{1565C0}     % dedicated color for .app file

% ============================================================
% Architecture Diagram
% ============================================================
\begin{document}

\begin{tikzpicture}[scale=0.68, every node/.style={scale=0.68},
    quadrant/.style={minimum width=6.5cm, minimum height=4.5cm, text width=5.8cm, align=left, font=\small},
    api/.style={font=\scriptsize\ttfamily, inner sep=2pt},
]
    % ---- Outer frames ----
    % Left half: User-defined code
    \node[draw=erlpurple!60, fill=erlpurple!3, rounded corners=6pt,
          minimum width=6.6cm, minimum height=10.4cm] (leftframe) at (-4.8, -0.55) {};
    \node[font=\small\bfseries\color{erlpurple}, anchor=south] at (leftframe.north) {用户定义模块（Application的启动）};
    \node[font=\small\ttfamily\bfseries\itshape\color{black}, anchor=north west] at ([xshift=3pt, yshift=-1pt]leftframe.north west) {bsc.erl + bsc.app};
    \draw[dashed, black, line width=0.5pt]
        ([xshift=3pt, yshift=-13pt]leftframe.north west) -- ([xshift=-3pt, yshift=-13pt]leftframe.north east);

    % Right half: OTP Application library
    \node[draw=erlblue!60, fill=erlblue!3, rounded corners=6pt,
          minimum width=6.6cm, minimum height=10.4cm] (rightframe) at (6.4, -0.55) {};
    \node[font=\small\bfseries\color{erlblue}, anchor=south] at (rightframe.north) {application库（应用管理引擎）};
    \node[font=\small\ttfamily\bfseries\itshape\color{black}, anchor=north west] at ([xshift=3pt, yshift=-1pt]rightframe.north west) {application.erl};
    \draw[dashed, black, line width=0.5pt]
        ([xshift=3pt, yshift=-13pt]rightframe.north west) -- ([xshift=-3pt, yshift=-13pt]rightframe.north east);

    % Connect the two dashed lines in the middle
    \draw[dashed, black, line width=0.5pt]
        ([xshift=-3pt, yshift=-13pt]leftframe.north east) -- ([xshift=3pt, yshift=-13pt]rightframe.north west);

    % ---- Left-Top: Callback Module + .app file (User Code — source) ----
    \node[draw=erlgreen, fill=erlgreen!8, rounded corners=5pt,
          minimum width=6.2cm, minimum height=4.2cm, text width=5.8cm, align=left] (lt) at (-4.8, 1.75) {};
    \node[font=\small\bfseries\color{erlgreen!80!black}, anchor=north west] at ([xshift=3pt, yshift=-3pt]lt.north west) {回调模块 + .app 资源文件};
    \node[font=\scriptsize\ttfamily, text width=5.4cm, align=left, anchor=north west] at ([xshift=6pt, yshift=-18pt]lt.north west) {
        \textcolor{erlgreen!70!black}{-module(bsc).}\\
        \textcolor{erlgreen!70!black}{-behaviour(application).}
    };
    % Highlight start/2 callback inside Callback Module
    \node[draw=erlred!70, fill=erlred!8, rounded corners=3pt,
          font=\scriptsize\ttfamily\bfseries, inner sep=3pt, anchor=north]
        (startcb) at ([yshift=-42pt]lt.north)
        {\textcolor{erlred!80!black}{\ding{72} start(\_Type, \_Args) -> bsc\_sup:start\_link()}};
    \node[font=\scriptsize\ttfamily, text width=5.4cm, align=left, anchor=north west] at ([xshift=6pt, yshift=-60pt]lt.north west) {
        \textcolor{erlinit}{\textbf{stop/1}}\\[2pt]
        {\tiny\textcolor{gray}{必须实现的回调函数：}}\\
        {\tiny\textcolor{gray}{start/2 → \{ok, Pid\} | \{ok, Pid, State\}}}\\
        {\tiny\textcolor{gray}{stop/1 → 清理工作（进程已终止后调用）}}
    };

    % Vertical dashed line from left endpoint of the dashed line down to just below lt.south
    \draw[dashed, black, line width=0.5pt]
        let \p1 = ([xshift=3pt, yshift=-13pt]leftframe.north west),
            \p2 = ([yshift=-6pt]lt.south)
        in (\p1) -- (\x1, \y2);

    % Horizontal dashed line from the bottom of the vertical line to near right outer frame left edge
    \draw[dashed, black, line width=0.5pt]
        let \p1 = ([xshift=3pt, yshift=-13pt]leftframe.north west),
            \p2 = ([yshift=-6pt]lt.south),
            \p3 = ([xshift=-3pt]rightframe.west)
        in (\x1, \y2) -- (\x3, \y2);

    % ---- Left-Bottom: Client API (runtime communication) ----
    \node[draw=erlpurple, fill=erlpurple!8, rounded corners=5pt,
          minimum width=6.2cm, minimum height=2.8cm, text width=5.8cm, align=left] (lb) at (-4.8, -2.85) {};
    \node[font=\small\bfseries\color{erlpurple}, anchor=north west] at ([xshift=3pt, yshift=-3pt]lb.north west) {Client进程（管理Application）};
    \node[font=\scriptsize\ttfamily, text width=5.4cm, align=left, anchor=north west] at ([xshift=6pt, yshift=-18pt]lb.north west) {
        application:start(\textcolor{erlapp}{\textbf{bsc}})\\
        application:stop(\textcolor{erlapp}{\textbf{bsc}})\\
        application:ensure\_all\_started(\textcolor{erlapp}{\textbf{bsc}})\\
        application:get\_env(\textcolor{erlapp}{\textbf{bsc}}, Key)
    };
    % Note: AppName = atom()
    \node[font=\tiny\color{gray}, anchor=south west] at ([xshift=6pt, yshift=3pt]lb.south west)
        {备注：\texttt{\textcolor{erlapp}{\textbf{AppName}}} = \texttt{atom()}，对应 .app 文件中的名称};

    % ---- Right-Top: Application Controller + Application Master ----
    \node[draw=erlblue, fill=erlblue!10, rounded corners=5pt,
          minimum width=6.2cm, minimum height=4.2cm, text width=5.8cm, align=left] (rt) at (6.4, 1.75) {};
    \node[font=\small\bfseries\color{erlblue}, anchor=north west] at ([xshift=3pt, yshift=-3pt]rt.north west) {Application Controller \& Master};
    \node[font=\scriptsize, text width=5.4cm, align=left, anchor=north west] at ([xshift=6pt, yshift=-16pt]rt.north west) {
        \textcolor{erlblue!80!black}{\ding{72} \textbf{Application Controller（单例）：}}\\
        {\tiny\textcolor{erlred!80!black}{\textbf{每个节点一个，管理所有 Application}}}\\
        {\tiny 加载 .app $\rightarrow$ 检查依赖 $\rightarrow$ 创建 App Master}\\[6pt]
        \textcolor{erlblue!80!black}{\ding{72} \textbf{Application Master（每App一个）：}}\\
        {\tiny 调用 \textcolor{erlmodule}{\textbf{\texttt{Module:start/2}}} $\rightarrow$ 获取 Top Supervisor Pid}\\
        {\tiny 作为 Group Leader，监控 Top Supervisor}\\[6pt]
        \textcolor{erlblue!80!black}{\ding{72} \textbf{启动流程：}}\\
        {\tiny\texttt{application:start} $\rightarrow$ \texttt{load .app} $\rightarrow$ \texttt{check deps}}\\
        {\tiny $\rightarrow$ \texttt{create Master} $\rightarrow$ \textcolor{erlmodule}{\textbf{\texttt{Module:start/2}}} $\rightarrow$ \texttt{\{ok, Pid\}}}
    };

    % ---- Right-Bottom: Supervision Tree (managed by Application) ----
    \node[draw=erlgreen!70!black, fill=erlgreen!5, rounded corners=5pt,
          minimum width=6.2cm, minimum height=2.8cm, text width=5.8cm, align=left] (rb) at (6.4, -2.85) {};
    \node[font=\small\bfseries\color{erlgreen!70!black}, anchor=north west] at ([xshift=3pt, yshift=-3pt]rb.north west) {监督树（由Application管理）};
    \node[font=\tiny, text width=5.4cm, align=left, anchor=north west] at ([xshift=6pt, yshift=-16pt]rb.north west) {
        \textcolor{erlred!80!black}{\texttt{Module:start/2}} 返回后：\\[3pt]
        \textcolor{erlblue}{\ding{72}} \texttt{Top Supervisor:} \textcolor{erlblue}{\textbf{bsc\_sup}}\\
        \hspace{8pt}\textcolor{erlgreen!70!black}{\ding{72}} \texttt{child\_1:} \textcolor{erlgreen!70!black}{\textbf{gen\_server}} $\rightarrow$ \texttt{frequency}\\
        \hspace{8pt}\textcolor{erlpurple}{\ding{72}} \texttt{child\_2:} \textcolor{erlpurple}{\textbf{gen\_event}} $\rightarrow$ \texttt{freq\_overload}\\
        \hspace{8pt}\textcolor{erlinit}{\ding{72}} \texttt{child\_3:} \textcolor{erlinit}{\textbf{gen\_statem}} $\rightarrow$ \texttt{phone\_fsm}\\[2pt]
        {\fontsize{4}{5}\selectfont\textcolor{gray}{{\fontsize{3}{4}\selectfont Top Supervisor} 终止 → {\fontsize{3}{4}\selectfont Application Master} 检测 → {\fontsize{3}{4}\selectfont Application} 终止}}
    };

    % ==== .app Resource File Info Box (center-bottom gap) ====
    \node[draw=erlyellow!80!black, fill=erlyellow!10, rounded corners=5pt,
          minimum width=4.0cm, minimum height=3.4cm, text width=3.6cm, align=left,
          font=\tiny] (appfilebox) at (0.8, -3.45) {
        \textcolor{erlyellow!70!black}{\scriptsize\bfseries .app 资源文件}\\[2pt]
        \texttt{\{application, bsc,}\\[0pt]
        \hspace{4pt}\texttt{[\{}\textcolor{erlblue}{\textbf{mod}}\texttt{, \{bsc, []\}\},}\\
        \hspace{5pt}\texttt{\{}\textcolor{erlred}{\textbf{applications}}\texttt{,}\\
        \hspace{10pt}\texttt{[kernel, stdlib, sasl]\},}\\
        \hspace{5pt}\texttt{\{}\textcolor{erlgreen}{\textbf{env}}\texttt{, [...]\},}\\
        \hspace{5pt}\texttt{\{}\textcolor{erlpurple}{\textbf{modules}}\texttt{, [bsc, ...]\},}\\
        \hspace{5pt}\texttt{\{}\textcolor{erlinit}{\textbf{registered}}\texttt{, [...]\}}\\
        \texttt{]\}.}\\[3pt]
        \ding{72} \textcolor{erlblue}{\textbf{mod}} 指定回调模块\\
        \ding{72} \textcolor{erlred}{\textbf{applications}} 声明依赖
    };

    % ==== Arrows ====

    % ==== Phase 1: Startup ====

    % Arrow ①: application:start — from Client (lb) to Application Controller (rt)
    \draw[-{Stealth[length=6pt]}, ultra thick, erlred!80!black, line width=1.6pt]
        ([yshift=3.8mm]lt.east) -- node[above, font=\scriptsize, align=center, text=erlred!80!black]
        {\textcolor{black}{\ding{172}} 启动 Application\\[1pt]
         \texttt{application:start(\textcolor{erlapp}{\textbf{bsc}})}} ([yshift=3.8mm]rt.west);

    % Arrow ③: Return {ok, Pid} — from Application Master back to Client
    \draw[-{Stealth[length=6pt]}, ultra thick, erlpid!90!black, line width=1.6pt]
        ([yshift=-6.5mm]rt.west) -- node[below, font=\scriptsize, align=center, text=erlpid!80!black]
        {\textcolor{black}{\ding{174}} 返回 \texttt{\textcolor{erlpid}{\textbf{\{ok, Pid\}}}}\\[1pt]
         {{\fontsize{4}{5}\selectfont \textcolor{erlmodule}{\textbf{Module:start/2}} 成功后，Application 进入运行状态}}} ([yshift=-6.5mm]lt.east);

    % Arrow ④: Client API (lb) -> Callback Module (lt) (manage application)
    \draw[-{Stealth[length=5pt]}, thick, erlblue, line width=1.2pt]
        ([xshift=-4mm]lb.north) -- node[left, font=\scriptsize, align=center, anchor=north east, yshift=6pt] {\textcolor{black}{\ding{175}} \textcolor{erlblue}{start / stop /}\\[1pt]\textcolor{erlblue}{get\_env}} ([xshift=-4mm]lt.south);

    % Arrow ②: Application Master -> Supervision Tree (start tree)
    \draw[-{Stealth[length=6pt]}, thick, erlgreen!70!black, line width=1.2pt]
        ([xshift=0mm]rt.south) -- node[left, font=\scriptsize, align=right, text=erlgreen!70!black]
        {{\scriptsize\bfseries \textcolor{black}{\ding{173}} 启动监督树}}
        node[right, font=\tiny, align=left, text=erlgreen!70!black]
        {Module:start/2\\$\rightarrow$ bsc\_sup:start\_link()} ([xshift=0mm]rb.north);

    % Arrow ⑤: Callback Module (lt) -> Client API (lb) (Reply)
    \draw[-{Stealth[length=5pt]}, thick, erlred, line width=1.2pt]
        ([xshift=4mm]lt.south) -- node[right, font=\scriptsize, align=center, anchor=north west, yshift=6pt] {\textcolor{black}{\ding{176}} \textcolor{erlred}{ok / \{error, R\}}} ([xshift=4mm]lb.north);

    % ---- Dashed separator lines below inner boxes ----
    \draw[dashed, erlpurple!40, line width=0.6pt]
        ([yshift=-6pt]lb.south west) -- ([yshift=-6pt]lb.south east);
    \draw[dashed, black, line width=0.6pt]
        let \p1 = ([xshift=-3pt]rightframe.west),
            \p2 = ([yshift=-6pt]rb.south west),
            \p3 = ([xshift=-3pt]rightframe.east)
        in (\x1, \y2) -- (\x3, \y2);

    % Vertical dashed line from the right endpoint of horizontal dashed line down to the right bottom dashed line
    \draw[dashed, black, line width=0.5pt]
        let \p1 = ([xshift=-3pt]rightframe.west),
            \p2 = ([yshift=-6pt]lt.south),
            \p3 = ([yshift=-6pt]rb.south)
        in (\x1, \y2) -- (\x1, \y3);

    % Vertical dashed line from top dashed line right endpoint down to right bottom dashed line right endpoint
    \draw[dashed, black, line width=0.5pt]
        let \p1 = ([xshift=-3pt, yshift=-13pt]rightframe.north east),
            \p2 = ([yshift=-6pt]rb.south),
            \p3 = ([xshift=-3pt]rightframe.east)
        in (\x3, \y1) -- (\x3, \y2);

    % Left bottom annotation (inside left outer frame, below dashed line)
    \node[font=\tiny\color{erlred!70!black}, anchor=north west, text width=5.8cm, align=left] at ([xshift=6pt, yshift=-10pt]lb.south west)
        {\ding{72} \texttt{.app} 文件中的 \texttt{\textcolor{erlblue}{\textbf{mod}}} 字段是连接左右两半的关键桥梁\\
         \ding{72} \texttt{\{mod, \{bsc, []\}\}} 告诉 Application Master 调用 \texttt{bsc:start/2}\\
         \ding{72} Application 只需实现 \texttt{start/2} 和 \texttt{stop/1} 两个回调\\
         \ding{72} 可选回调：\texttt{prep\_stop/1}、\texttt{start\_phase/3}、\texttt{config\_change/3}};

    % Right bottom annotation (inside right outer frame, below dashed line)
    \node[font=\fontsize{5}{6}\selectfont\color{erlred!70!black}, anchor=north west, text width=5.8cm, align=left] at ([xshift=6pt, yshift=-10pt]rb.south west)
        {\ding{72} {\fontsize{4}{5}\selectfont Application Controller} 加载 \texttt{.app} 文件，检查 \texttt{applications} 依赖\\
         \ding{72} {\fontsize{4}{5}\selectfont Application Master} 调用 \textcolor{erlmodule}{\textbf{\texttt{Module:start/2}}} 启动监督树\\
         \ding{72} 停止时：\texttt{prep\_stop/1} $\rightarrow$ 关闭监督树 $\rightarrow$ \texttt{stop/1}\\
         \ding{72} restart\_type: \texttt{permanent} $|$ \texttt{transient} $|$ \texttt{temporary}};
\end{tikzpicture}

\end{document}
